\section{USAMO 1983}
        \begin{enumerate}
        	\item (\textit{USAMO 1983}) On a given circle, six points $A$, $B$, $C$, $D$, $E$, and $F$ are chosen at random, independently and uniformly with respect to arc length. Determine the probability that the two triangles $ABC$ and $DEF$ are disjoint, i.e., have no common points.


 



\textbf{Solution}

First we give the circle an orientation (e.g., letting the circle be the unit circle in polar coordinates).  Then, for any set of six points chosen on the circle, there are exactly $6!$ ways to label them one through six. Also, this does not affect the probability we wish to calculate. This will, however, make calculations easier.

 Note that, for any unordered set of six points chosen from the circle boundary, the number of ways to number them such that they satisfy this disjoint-triangle property is constant: there are six ways to choose which triangle will be numbered with the numbers one through three, and there are $(3!)^2$ ways to arrange the numbers one through three and four through six on these two triangles. Therefore, for any given configuration of points, there are $6^3=216$ ways to label them to have this disjoint-triangle property. There are, however, $6!=720$ ways to label the points in all, so given any six unordered points, the probability that when we inflict an ordering we produce the disjoint-triangle property is $216/720=3/10$.

 Since this probability is constant for any configuration of six unordered points we choose, we must have that $3/10$ is the probability that we produce the disjoint-triangle property if we choose the points as detailed in the problem statement.

 


	\item (\textit{USAMO 1983}) Prove that the roots of 
 \[x^5 + ax^4 + bx^3 + cx^2 + dx + e = 0\] cannot all be real if $2a^2 < 5b$.


 



\textbf{Solution}

We prove the contrapositive: if the polynomial in question has the five real roots $x_1, x_2, x_3, x_4, x_5$, then $5b \le 2a^2$.

 Because $a = -(x_1 + x_2 + x_3 + x_4 + x_5)$ and $b = x_1x_2 + x_1x_3 + ... + x_4x_5$ by Vieta's Formulae, we have

 
 \[2b = 2x_1x_2 + 2x_1x_3 + ... + 2x_4x_5 = (x_1 + x_2 + x_3 + x_4 + x_5)^2 - (x_1^2 + x_2^2 + x_3^2 + x_4^2 + x_5^2)\]
 \[=a^2 - \frac{(1+1+1+1+1)(x_1^2 + x_2^2 + x_3^2 + x_4^2 + x_5^2)}{5}\]
 \[\le a^2 - \frac{(x_1 + x_2 + x_3 + x_4 + x_5)^2}{5}\] (by Cauchy-Schwarz)
 \[=\frac{4a^2}{5},\]

 so $5b \le 2a^2$, as desired.

 



\textbf{Solution 2}

\textbf{Lemma:}

 For all real numbers $x_1,x_2,\cdots x_5$,

 
 \[2(x_1^2+x_2^2+\cdots+x_5^2)\ge\]

 
 \[x_1x_2+x_1x_3+\cdots+x_4x_5\]

 By the trivial inequality,

 
 \[x^2+y^2\ge 2xy \Rightarrow \frac{x^2}{2} + \frac{y^2}{2} \ge xy\]

 Making such an inequality for all the variable pairs and summing them, we find the lemma is true.

 Now, let our roots be $x_1,x_2,\cdots,x_5$. By Vieta's, $a=x_1+x_2+\cdots+x_5$ and $b=x_1x_2+x_1x_3+\cdots+x_4x_5$

 If we show that for all real $x_1,x_2,\cdots, x_5$ that $2a^2\ge 5b$, then we have a contradiction and all of $x_1,x_2,\cdots, x_5$ cannot be real. We start by rewriting $2a^2\ge 5b$ as

 
 \[2(x_1+x_2+\cdots+x_5)^2\ge 5(x_1x_2+x_1x_3+\cdots+x_4x_5)\]

 We divide by $2$ and find

 
 \[(x_1+x_2+\cdots+x_5)^2\ge \frac{5}{2}(x_1x_2+x_1x_3+\cdots+x_4x_5)\]

 Expanding the LHS, we have

 
 \[x_1^2+x_2^2+\cdots+x_5^2+2(x_1x_2+x_1x_3+\cdots+x_4x_5)\ge\frac{5}{2}(x_1x_2+x_1x_3+\cdots+x_4x_5)\]

 We subtract the sum in brackets, and then multiply by $2$ to find

 
 \[2x_1^2+2x_2^2+\cdots+2x_5^2\ge x_1x_2+x_1x_3+\cdots+x_4x_5\]

 which is true by our lemma.

 


	\item (\textit{USAMO 1983}) Each set of a finite family of subsets of a line is a union of two closed intervals. Moreover, any three of the sets of the family have a point in common. Prove that there is a point which is common to at least half the sets of the family.


 



\textbf{Solution}

Let us first see that this works for anything less than three sets.

 Obviously, due to the second condition given, they all share at least one point in common.

 Now, let us see that it works for three or more sets.

 First, take any three of the subsets in the family of sets.

 Let them be A, B, and C. We know that they share at least one point in common. Let us take the pair of segments that have the shortest length when you take the union of the pair. Let us denote this segment D. Evidently, if you take any other set of points, it must intersect D at least at one point. In order to minimize the number of intersections between the other varying segments, we must have at close to 1/2 of the remaining segments as possible intersecting each other. However, if we note, putting them that way forces at least 1/2 of the segments in each section.

 


	\item (\textit{USAMO 1983}) Six segments $S_1, S_2, S_3, S_4, S_5,$ and $S_6$ are given in a plane. These are congruent to the edges $AB, AC, AD, BC, BD,$ and $CD$, respectively, of a tetrahedron $ABCD$. Show how to construct a segment congruent to the altitude of the tetrahedron from vertex $A$ with straight-edge and compasses.


 



\textbf{Solution}

Throughout this solution, we denote the length of a segment $S$ by $|S|$.

 In this solution, we employ several lemmas. Two we shall take for granted: given any point $A$ and a line $\ell$ not passing through $A$, we can construct a line $\ell'$ through $A$ parallel to $\ell$; and given any point $A$ on a line $\ell$, we can construct a line $\ell'$ through $A$ perpendicular to $\ell$.

 \textbf{Lemma 1:} If we have two segments $S$ and $T$ on the plane with non-zero length, we may construct a circle at either endpoint of $S$ whose radius is $|T|$.

 \textbf{Proof:} We can construct arbitrarily many copies of $T$ by drawing a circle about one of its endpoints through its other endpoint, and then connecting the center of this circle with any other point on the circle. We can construct copies of $T$ like this until we create a circle of radius $|T|$ and center $P_1$ that intersects segment $S$. We can then take this intersection point $P_2$ and draw a line $\ell$ through it perpendicular to $S$, and draw a circle with center $P_2$ passing through $P_1$, and consider its intersection $P_3$ with $\ell$. Note that $P_2P_3\perp S$ and $|P_2P_3|=|T|$. Take an endpoint $P_4$ of $S$: then draw a line through $P_3$ parallel to $S$, and a line through $P_4$ parallel to $P_2P_3$. Let these two lines intersect at $P_5$. Then $P_2P_3P_5P_4$ is a rectangle, so $|P_4P_5|=|T|$. Our desired circle is then a circle centered at $P_4$ through $P_5$.

 \textbf{Lemma 2:} Given three collinear points $A$, $B$, $C$ in this order, if $|AB|=a$ and $|BC|=b$ with $a>b$, then we can construct a segment of length $\sqrt{a^2-b^2}$.

 \textbf{Proof:} From Lemma 1, we can construct a circle through $C$ with radius $a$, and then construct a perpendicular through $B$ to $AC$: these two objects intersect at $D$ and $E$. Both $BD$ and $DE$ have length $\sqrt{a^2-b^2}$, from the Pythagorean Theorem.

 \textbf{Proof of the original statement:} Note that we can construct a triangle $B'C'D'$ congruent to triangle $BCD$ by applying Lemma 1 to segments $S_4$, $S_5$, and $S_6$. Similarly, we can construct $A_B$ and $A_C$ outside triangle $B'C'D'$ such that $A_BC'D'\cong ACD$ and $A_CB'D'\cong ABD$.

 Let $A'$ be a point outside of the plane containing $S_1$ through $S_6$ such that $A'B'C'D'\cong ABCD$. Then the altitudes of triangles $A_BC'D'$ and $A'C'D'$ to segment $C'D'$ are congruent, as are the altitudes of triangles $A_CB'D'$ and $A'B'D'$ to segment $B'D'$. However, if we project the altitudes of $A'C'D'$ and $A'B'D'$ from $A'$ onto the plane, their intersection is the base of the altitude of tetrahedron $A'B'C'D'$ from $A'$. In addition, these altitude projections are collinear with the altitudes of triangles $A_BC'D'$ and $A_CB'D'$. Therefore, the altitudes of $A_BC'D'$ and $A_CB'D'$ from $A_B$ and $A_C$ intersect at the base $X'$ of the altitude of $A'B'C'D'$ from $A'$. In summary, we can construct $X'$ by constructing the perpendiculars from $A_B$ and $A_C$ to $C'D'$ and $B'D'$ respectively, and taking their intersection.

 Let $Y'$ be the intersection of $A_BX'$ with $C'D'$. Then the altitude length we seek to construct is, from the Pythagorean Theorem, $\sqrt{|A_BY'|^2-|X'Y'|^2}$. We can directly apply Lemma 2 to segment $A_BY'X'$ to obtain this segment. This shows how to construct a segment of length $A'X'$.

 


	\item (\textit{USAMO 1983}) Consider an open interval of length $1/n$ on the real number line, where $n$ is a positive integer. Prove that the number of irreducible fractions $p/q$, with $1\le q\le n$, contained in the given interval is at most $(n+1)/2$.


 



\textbf{Solution}

Let $I$ be an open interval of length $1/n$ and $F_n$ the set of fractions $p/q\in I$ with $p,q\in\mathbb{Z}$, $\gcd(p,q)=1$ and $1\leq q\leq n$. 

 Assume that $\frac{p}{q}\in F_n$. If $k\in\mathbb{Z}$ is such that $1\leq kq\leq n$, and $p'\in\mathbb{Z}$ is such that $\gcd(p',kq)=1$, then
 \[\left|\frac{p}{q}-\frac{p'}{kq}\right|\geq\frac{1}{kq}\geq \frac{1}{n}\]Therefore $\frac{p'}{kq}\notin I\supset F_n$. This means that $\frac{p}{q}$ is the only fraction in $F_n$ with denominator $q$ or multiple of $q$.

 Therefore, from each of the pairs in $P=\left\{(k,2k):\ 1\leq k\leq \left\lfloor\frac{n+1}{2}\right\rfloor\right\}$ at most one element from each can be a denominator of a fraction in $F_n$.

 Hence $|F_n|\leq |P|\leq\frac{n+1}{2}$

 


\end{enumerate}